% ============================================================
%  List motywacyjny — Maciej Górski
%  Styl dopasowany do cv-maciej-gorski.tex
%  Kompilacja: XeLaTeX lub LuaLaTeX (zalecane)
%  Overleaf:   Menu → Compiler → XeLaTeX
%
%  UWAGA: Jeśli nie masz fontu Inter, użyj Noto Sans:
%    \setmainfont{Noto Sans}
% ============================================================
\documentclass[11pt,a4paper]{article}

% ---------- pakiety ----------
\usepackage[polish]{babel}
\usepackage{fontspec}
\setmainfont{Inter}[
  Extension      = .ttf,
  UprightFont    = *-Regular,
  BoldFont       = *-Bold,
  ItalicFont     = *-Italic,
  BoldItalicFont = *-BoldItalic,
]
\usepackage{fontawesome5}
\usepackage{xcolor}
\usepackage{hyperref}
\usepackage{tcolorbox}
\tcbuselibrary{skins}
\usepackage[scale=0.85]{geometry}
\usepackage{parskip}
\usepackage{setspace}
\onehalfspacing

% ---------- kolory (złoty akcent) ----------
\definecolor{sidebar}{HTML}{0a0f1c}
\definecolor{sidebarText}{HTML}{e2e8f0}
\definecolor{goldAccent}{HTML}{e5b05c}
\definecolor{goldDark}{HTML}{c9943a}
\definecolor{mutedText}{HTML}{94a3b8}

% ---------- hyperlinki ----------
\hypersetup{colorlinks=true, urlcolor=goldAccent, linkcolor=goldAccent}

% ---------- dokument ----------
\begin{document}
\thispagestyle{empty}

% ======================== PASEK BOCZNY ========================
\begin{tcolorbox}[
  colback      = sidebar,
  coltext      = sidebarText,
  boxrule      = 0pt,
  leftrule     = 0pt,
  arc          = 0pt,
  outer arc    = 0pt,
  width        = 0.28\paperwidth,
  left         = 0.8em,
  right        = 0.8em,
  top          = 2em,
  bottom       = 2em,
  after        = \hspace{0.5em},
  nobeforeafter,
  enhanced,
  remember,
]
% ---- inicjały ----
\vspace*{-0.5em}
\begin{center}
  \fcolorbox{goldAccent!60}{sidebar!90}{%
    \parbox{3.2cm}{\centering\vspace{1.2cm}
      {\fontsize{48}{48}\selectfont\bfseries\color{goldAccent} MG}\vspace{1.2cm}}
  }
\end{center}

\vspace{0.6em}

% ---- kontakt ----
{\small\bfseries\color{goldAccent} \faIcon[regular]{envelope}\hspace{0.4em}KONTAKT}\\[0.25em]
{\footnotesize
\faIcon[regular]{envelope}\hspace{0.4em}\href{mailto:maciejgorski978@gmail.com}{maciejgorski978@gmail.com}\\[0.15em]
\faIcon{phone-alt}\hspace{0.4em}+48 514 299 431\\[0.15em]
\faIcon{map-marker-alt}\hspace{0.4em}Polska
}

\vspace{0.8em}

% ---- online ----
{\small\bfseries\color{goldAccent} \faIcon{globe}\hspace{0.4em}ONLINE}\\[0.25em]
{\footnotesize
\faIcon{github}\hspace{0.4em}\href{https://github.com/Gorski-Maciej}{github.com/Gorski-Maciej}\\[0.15em]
\faIcon{linkedin}\hspace{0.4em}\href{https://linkedin.com/in/maciej-gorski-046176282}{LinkedIn}
}

\vspace{0.8em}

% ---- języki ----
{\small\bfseries\color{goldAccent} \faIcon{language}\hspace{0.4em}JĘZYKI}\\[0.25em]
{\footnotesize
\textbf{Polski} --- ojczysty\\[0.1em]
\textbf{Angielski} --- B1/B2
}

\vspace{0.8em}

% ---- kluczowe technologie ----
{\small\bfseries\color{goldAccent} \faIcon{code}\hspace{0.4em}TECHNOLOGIE}\\[0.25em]
{\footnotesize
\setlength{\tabcolsep}{0.35em}
\begin{tabular}{ll}
\color{goldAccent}\textbullet & Python / FastAPI / Flask \\
\color{goldAccent}\textbullet & Docker / Kubernetes \\
\color{goldAccent}\textbullet & CI/CD / GitHub Actions \\
\color{goldAccent}\textbullet & PostgreSQL / Redis / Kafka \\
\color{goldAccent}\textbullet & Linux / Bash \\
\color{goldAccent}\textbullet & React / Vite \\
\color{goldAccent}\textbullet & AWS \\
\end{tabular}
}

\vspace{0.8em}

% ---- sieci ----
{\small\bfseries\color{goldAccent} \faIcon{network-wired}\hspace{0.4em}SIEĆ}\\[0.25em]
{\footnotesize
\setlength{\tabcolsep}{0.35em}
\begin{tabular}{ll}
\color{goldAccent}\textbullet & TCP/IP, VLAN, OSPF \\
\color{goldAccent}\textbullet & Cisco / GNS3 \\
\color{goldAccent}\textbullet & Firewall / Wireshark \\
\color{goldAccent}\textbullet & CCNA (w trakcie) \\
\end{tabular}
}

\vspace{0.8em}

% ---- kompetencje ----
{\small\bfseries\color{goldAccent} \faIcon{users}\hspace{0.4em}KOMPETENCJE}\\[0.25em]
{\footnotesize
Analityczne myślenie\\[0.05em]
Samodzielność\\[0.05em]
Szybka nauka\\[0.05em]
Odpowiedzialność\\[0.05em]
Praca zespołowa
}
\end{tcolorbox}
%
\begin{minipage}[t]{0.68\paperwidth}
\vspace{1.2em}
\hspace{1.5em}
\begin{minipage}[t]{0.9\textwidth}

% ======================== NAGŁÓWEK ========================
{\fontsize{28}{32}\selectfont\bfseries\color{sidebar} Maciej Górski\par}
\vspace{0.15em}
{\large\color{goldAccent!80!black} Student Informatyki \textbar\ Python Developer \textbar\ Network Enthusiast\par}

\vspace{1.2em}

% ---- miejsce i data ----
{\small\color{mutedText}
\faIcon{map-marker-alt}\hspace{0.3em}[Miejscowość], \today
}

\vspace{1.5em}

% ---- dane adresata (do uzupełnienia) ----
{\small
\textbf{\faIcon{building}\hspace{0.3em}Dane adresata:}\\[0.3em]
[Nazwa firmy]\\[0.1em]
[Adres]\\[0.1em]
[Kod pocztowy, Miasto]
}

\vspace{1.5em}

% ---- tytuł ----
\noindent{\large\bfseries\color{sidebar} List motywacyjny}

\vspace{0.4em}
\noindent\textbf{Dotyczy:} odpowiedź na ogłoszenie o staż / praktykę w obszarze IT

\vspace{0.8em}

% ---- treść ----
Szanowni Państwo,\\[0.6em]

z dużą satysfakcją odpowiadam na ogłoszenie o staż/praktykę w Państwa firmie. Jestem studentem pierwszego roku Informatyki (specjalizacja: administracja sieci i systemów) i poszukuję możliwości rozwoju w obszarach \textbf{Python, DevOps, administracja Linux} oraz \textbf{sieci komputerowe}. Państwa firma, jako lider w swojej branży, jest miejscem, w którym mógłbym łączyć swoją pasję do programowania z zainteresowaniem infrastrukturą IT.

\vspace{0.4em}

W ramach własnych projektów zrealizowałem cztery kompleksowe systemy mikrousługowe w architekturze Docker Compose:

\begin{itemize}
  \item \textbf{CloudBudget} --- platforma FinOps z FastAPI, PostgreSQL, Celery i React. System analizuje koszty chmurowe, generuje rekomendacje ML i udostępnia dashboard operatorski.
  \item \textbf{NetGuardian} --- system SecOps/SIEM/SOAR wykorzystujący Kafka do ingestu zdarzeń, Redis do cache'owania, DuckDB do analizy i ML (scikit-learn) do detekcji anomalii. Zawiera enrichment GeoIP oraz automatyczne playbooki reakcji.
  \item \textbf{NetAegis} --- dwuwarstwowa orkiestracja agentów w architekturze MCP (Model Context Protocol) z frontendem React i komunikacją WebSocket w czasie rzeczywistym.
  \item \textbf{InfraFlow} --- platforma monitorowania infrastruktury z metrykami Prometheus i dashboardami Grafana.
\end{itemize}

Projekty te realizowałem samodzielnie, dbając o testy (pytest), dokumentację i ciągłą integrację (GitHub Actions). Każdy z nich pokazuje moją umiejętność pracy z nowoczesnym stackiem technologicznym: \textbf{FastAPI, Docker, PostgreSQL, Redis, Kafka, Prometheus, Grafana, React, CI/CD} oraz chmurą \textbf{AWS}.

\vspace{0.4em}

Równolegle rozwijam wiedzę sieciową --- przygotowuję się do certyfikacji \textbf{Cisco CCNA}, pracuję w GNS3 i Packet Tracer (VLAN, OSPF, ACL), a na co dzień używam Linuxa (Debian/Ubuntu), Bash i Git. Posiadam również praktyczne umiejętności z zakresu IoT, bezpieczeństwa IT oraz testowania oprogramowania.

\vspace{0.4em}

Cechuje mnie \textbf{analityczne myślenie, samodzielność i szybka nauka}. Potrafię efektywnie diagnozować problemy, dokumentować rozwiązania i pracować zarówno samodzielnie, jak i w zespole. Jestem otwarty na pracę hybrydową i zdalną.

\vspace{0.4em}

Chciałbym rozwijać się w Państwa firmie jako \textbf{Python Developer, specjalista IT Support / DevOps lub administrator sieci}. Gwarantuję pełne zaangażowanie, motywację do nauki i rzetelność w wykonywaniu powierzonych zadań.

\vspace{0.4em}

Będę wdzięczny za możliwość spotkania i rozmowy o moim potencjalnym udziale w projektach Państwa firmy.

\vspace{1em}

Z wyrazami szacunku,\\[0.4em]
{\bfseries\color{sidebar} Maciej Górski}

\vspace{0.2em}
{\footnotesize\color{mutedText}
\faIcon[regular]{envelope}\hspace{0.3em}\href{mailto:maciejgorski978@gmail.com}{maciejgorski978@gmail.com} \hspace{1em}
\faIcon{phone-alt}\hspace{0.3em}+48 514 299 431 \hspace{1em}
\faIcon{github}\hspace{0.3em}\href{https://github.com/Gorski-Maciej}{github.com/Gorski-Maciej}
}

\end{minipage}
\end{minipage}

\end{document}
